% !TEX root = template.tex

\clearpage
\section{Codebase and Execution Guide}
\label{sec:usage}

The project repository is structured to provide a robust, unified interface for local and remote execution via a central \textbf{\texttt{Makefile}}. The codebase supports three distinct paradigms for running the training and testing pipelines, seamlessly adapting the user's configurations across execution environments.

\subsection{Execution Environments}
\begin{enumerate}
    \item \textbf{Local Bare Metal:} Direct execution natively in the terminal using the local virtual environment (e.g., \textbf{\texttt{make train}}). Ideal for quick debugging and syntax checks.
    \item \textbf{Local Docker:} Wraps the execution within an isolated local container (e.g., \textbf{\texttt{make train-local}}). Ensures dependency consistency before remote deployment.
    \item \textbf{Hyperstack Cloud Docker:} Uses the \textbf{\texttt{remote\_exec/deploy.py}} script to securely connect, synchronize the local codebase to the cloud via \textbf{\texttt{rsync}}, and dynamically provision or resume a GPU cloud instance via the Hyperstack API before interactively executing training inside a container while streaming the logs back (e.g., \textbf{\texttt{make train-hyperstack}}).
\end{enumerate}

\subsection{Repository Structure and Dataset Placement}
The repository is modularly organized into the following key directories to separate the model logic from the execution environment:
\begin{itemize}
    \item \textbf{\texttt{architectures/}}: Contains the PyTorch definitions for the neural network models (e.g., \texttt{resnet.py}, \texttt{inception.py}, \texttt{transformer.py}).
    \item \textbf{\texttt{core/}}: Houses the core functionalities, including the custom dataset classes, data augmentations, and evaluation/training loops (e.g., \texttt{dopplerdataset.py}, \texttt{evaluation.py}).
    \item \textbf{\texttt{dataset/}}: The designated directory for the Doppler spectrogram dataset. \textbf{Important:} To ensure the data loaders function correctly, the raw dataset must be extracted and placed exactly at \texttt{dataset/data/doppler\_traces/}.
    \item \textbf{\texttt{remote\_exec/}}: Contains the Python deployment and synchronization scripts responsible for interacting with the Hyperstack API and managing the cloud instances.
    \item \textbf{\texttt{scripts/}}: Provides various utility scripts for integrating with Weights \& Biases, generating graphs, and performing system checks.
    \item \textbf{\texttt{weights/}}: The destination directory where serialized PyTorch model checkpoints (\texttt{.pt} files) are automatically saved during training and loaded from during testing.
    \item \textbf{\texttt{wandb/}}: Locally stores the synchronization logs for Weights \& Biases.
\end{itemize}

\subsection{Key Scripts and Configurations}
The fundamental logic resides in independent, modular scripts, which the \texttt{Makefile} coordinates:
\begin{itemize}
    \item \texttt{train.py}: The primary training pipeline. It dynamically parses arguments for architecture (\texttt{resnet8}, \texttt{inception}, \texttt{transformer}), task (\texttt{activity}, \texttt{person\_id}), learning rate, and epochs.
    \item \texttt{test.py}: Evaluates stored weights against the isolated test dataset, producing analytical confusion matrices and accuracy metrics.
    \item \texttt{check\_wandb.py} \& \texttt{download\_wandb\_history.py}: Integration scripts that query the Weights \& Biases API to retrieve performance history and compile comparative results and graphs.
\end{itemize}

All Makefile commands accept environmental variable overrides for rapid prototyping. For instance, testing a Transformer on the Activity Recognition task locally can be launched as simply as:

\vspace{2mm}
\noindent\fbox{%
    \parbox{\linewidth}{%
        \centering\textbf{\texttt{make train ARCH=transformer TASK=activity EPOCHS=10 LR=5e-4}}
    }%
}
\vspace{2mm}

\subsection{Remote Sync and Cleanup}
Remote operations automatically synchronize the latest codebase changes prior to execution. Following execution, results and serialized weights (\texttt{.pt} files) are retrieved using the provided fetch commands:
\begin{itemize}
    \item \textbf{\texttt{make pull-hyperstack}}: Retrieves all trained weights from the cloud.
    \item \textbf{\texttt{make pull-results-hyperstack}}: Retrieves just the test logs, CSVs, and generated graphs.
    \item \textbf{\texttt{make clean-all}}: Safely tears down local containers, remote containers, and halts cloud instances to prevent inactive billing.
\end{itemize}
